\section{Explicit-regularization baselines}
\label{sec:regularization}

The following results preserve the initial shrinkage investigation as baselines
for the optimization study. For fixed positive vectors $u,v$, consider the
original-scale estimator
\begin{equation}
 \label{eq:objective}
 \hat\beta = \arg\min_{\beta \in \R^p} \left\{ \frac12\beta^\top C\beta - z^\top\beta + t\sum_j u_j|\beta_j| + \frac r2\sum_j v_j\beta_j^2 \right\}, \qquad t,r
           \ge 0.
\end{equation}
Dividing centered feature $j$ by $s_j$ and transforming the fitted coefficient
back gives $u_j = s_j$ and $v_j = s_j^2$. This equivalence requires no
orthogonality. Allowing separate weights also covers penalty schemes such as
$u_j = v_j = C_{jj}$.

Weights and penalties are fixed with respect to $y$ throughout the distribution
calculations. Cross-validation and response-dependent penalty paths require
separate analysis. Write $\soft_k(x) = \sign(x)(|x| - k)_+$ for
soft-thresholding.

\begin{proposition}[Ridge in arbitrary dimension]
 \label{prop:ridge}
 Let $V = \diag(v)$ and $t = 0$. Then
 \begin{align*}
  \E\hat\beta           & = (C + rV)^{-1}C\beta^*,             \\
  \E\hat\beta - \beta^* & = -r(C + rV)^{-1}V\beta^*,           \\
  \Cov(\hat\beta)       & = \tau^2(C + rV)^{-1}C(C + rV)^{-1}.
 \end{align*}
 The mean and covariance formulas require only zero-mean errors with covariance
 $\sigma_\varepsilon^2 I$, rather than Gaussian errors.
\end{proposition}
\begin{proof}
 Differentiate the quadratic objective to obtain $\hat\beta = (C + rV)^{-1}z$,
 then take its mean and covariance.
\end{proof}
For $r > 0$ and $0 < k < 1$, the identity $(C + rV)^{-1}C = kI$ implies
$V = (1 - k)C/(kr)$. Thus a diagonal ridge penalty can give the same shrinkage
factor for every coefficient vector only when $C$ is diagonal.
